\documentclass[aspectratio=169, 10pt]{beamer}
% Shared Beamer setup for the three workshop decks.
% Build with: python tasks.py slides w1   (or compile slides/wN.tex from the slides folder)
%
% Build variants (set by tasks.py through latexmk -usepretex):
%   \finalbuild  hides PREP markers; tasks.py refuses a final build while any remain
%   \withnotes   shows speaker notes on a second screen
%   handout      beamer class option for the participant copy

\usetheme[progressbar=frametitle, numbering=fraction, block=fill]{metropolis}
\usepackage[T1]{fontenc}
\usepackage{booktabs, array, tabularx, pifont, xcolor, tikz, pgfplots}
\pgfplotsset{compat=1.18}
\usetikzlibrary{arrows.meta, positioning, fit, backgrounds, calc}
\usepackage[backend=biber, style=numeric-comp, sorting=none, maxbibnames=3, giveninits=true]{biblatex}
\addbibresource{../../../bibliography/references.bib}
\renewcommand*{\bibfont}{\tiny}
\setbeamertemplate{bibliography item}{\insertbiblabel}

\definecolor{okgreen}{HTML}{2E7D32}
\definecolor{badred}{HTML}{C62828}
\definecolor{warnamber}{HTML}{B26A00}
\definecolor{tierAcol}{HTML}{1B5E20}
\definecolor{tierBcol}{HTML}{0D47A1}
\definecolor{tierCcol}{HTML}{6A1B9A}
\definecolor{orgrecycler}{HTML}{00796B}
\definecolor{orglab}{HTML}{1565C0}
\definecolor{orgmaker}{HTML}{E65100}
\definecolor{orgbroker}{HTML}{C62828}
\definecolor{orgregulator}{HTML}{5D4037}
\setbeamercolor{alerted text}{fg=badred}

% Table marks used by generated results
\newcommand{\cmark}{\textcolor{okgreen}{\ding{51}}}
\newcommand{\xmark}{\textcolor{badred}{\ding{55}}}
\newcommand{\pmark}{\textcolor{warnamber}{\textbf{$\sim$}}}

% Evidence tiers (docs/evidence-policy.md): \ev{A}{key} prints a badge and a citation
\newcommand{\tier}[1]{\tikz[baseline=(t.base)]\node[fill=tier#1col, text=white, font=\tiny\bfseries,
  inner sep=1.1pt, rounded corners=1pt] (t) {#1};}
\newcommand{\ev}[2]{\tier{#1}\,\cite{#2}}
\newcommand{\pre}{{\tiny\itshape(preprint)}}
\newcommand{\grey}[1]{\tier{C}\,\cite{#1}\,{\tiny\itshape motivation and vocabulary only}}

% Preparation markers: visible while drafting, must be resolved before the final build
\newcommand{\prep}[1]{\ifdefined\finalbuild\else{\setlength{\fboxsep}{1.5pt}%
  \colorbox{warnamber!20}{\parbox{0.92\linewidth}{\tiny\textbf{PREP:} #1}}}\fi}

\ifdefined\withnotes
  \setbeameroption{show notes on second screen=right}
\fi

% Generated experiment output, with a visible fallback when it has not been run
\newcommand{\result}[1]{\IfFileExists{../results/#1.tex}{\input{../results/#1.tex}}%
  {\colorbox{badred!15}{\scriptsize Missing result \texttt{\detokenize{#1}}: run \texttt{python tasks.py run}}}}

\newcommand{\discussion}[3]{% {number}{question}{readings}
  \begin{frame}{Discussion #1}
    \begin{alertblock}{Question for the room}\large #2\end{alertblock}
    \vspace{2mm}{\small\textbf{Readings.} #3}
  \end{frame}}

\author{Aditya Sissodiya}
\institute{Lule\aa{} University of Technology \quad\textperiodcentered\quad Agentic AI Security, Identity, and Access Control (7.5 ECTS)}


\title{Where authority leaks}
\subtitle{Workshop 1 of 3 \textperiodcentered{} Agents, protocols, and identity}
\date{\prep{set the workshop date}}

\begin{document}

\maketitle

\begin{frame}{Today}
  \begin{columns}[T]
    \column{0.57\textwidth}
    \small
    \begin{tabular}{@{}l l@{}}
      0:00 & The question \\
      0:10 & Talk and live demos (40 min) \\
      0:50 & Hands-on: build the delegation chain \\
      1:20 & Break \\
      1:30 & Discussion 1: where should governance live? \\
      2:00 & Discussion 2: workload identity or provenance? \\
      2:30 & Examiner questions \\
    \end{tabular}
    \column{0.39\textwidth}
    \begin{block}{What you need}
      \small A laptop with Python 3.10+ and a clone of the course repository.
      Everything runs on the standard library:\\[2pt]
      {\footnotesize\texttt{python tasks.py run w1}}
    \end{block}
  \end{columns}
  \note{Start the projector demo terminal before the session. Check that results/ is fresh.}
\end{frame}

\begin{frame}{Learning outcomes for this workshop}
  \small After this workshop I should be able to show that I can:
  \begin{enumerate}
    \item[K1.1] locate the authority boundaries (principal, actor, resource, organization) in a cross-organizational agent workflow, and say when multi-agent is justified;
    \item[K1.2] describe the MCP authorization flow and the A2A Agent Card model from the normative text;
    \item[K1.3] distinguish workload identity, on-behalf-of authority, impersonation, and delegation;
    \item[S1.4] use a simulation to show where authority leaks in a standards-compliant flow;
    \item[S1.5] list, with sources, the authorization decisions MCP and A2A leave to implementers;
    \item[J1.6] argue where cross-organizational governance should live, and defend it.
  \end{enumerate}
\end{frame}

\section{The question}

\begin{frame}{When an agent acts, who is authorized?}
  \begin{columns}[T]
    \column{0.48\textwidth}
    \textbf{Classical access control}
    \begin{itemize}\small
      \item A principal logs in; a process acts with that principal's rights.
      \item Authority is fixed at session start and checked by one reference monitor.
    \end{itemize}
    \column{0.48\textwidth}
    \textbf{Agentic systems}
    \begin{itemize}\small
      \item Every tool call and hand-off passes authority to another component.
      \item Chains are composed at run time, often across organizations.
      \item The component that acts may follow instructions it found in data.
    \end{itemize}
  \end{columns}
  \vspace{4mm}
  \begin{alertblock}{Today's claim}
    The protocols standardize how to \emph{call} an agent or tool. They leave most of
    \emph{who may act for whom} to whoever deploys them.
  \end{alertblock}
  \note{Industry surveys are sometimes quoted here. Do not quote the IBM/Ponemon figures without the full methodology caveat (vendor-sponsored, self-reported convenience panel, construct not operationalized).}
\end{frame}

\begin{frame}{Vocabulary for all three workshops}
  \small
  \begin{tabularx}{\textwidth}{@{}l X@{}}
    \toprule
    \textbf{Principal} & The party whose authority is being exercised (Alice). \\
    \textbf{Actor} & The component actually making the request (an agent). \\
    \textbf{Resource} & What is accessed: here, fields of a Digital Product Passport. \\
    \textbf{Delegation} & Passing some authority to an actor, who acts \emph{as itself, for} the principal. \\
    \textbf{Impersonation} & The actor acts \emph{as} the principal; the actor disappears from view. \\
    \textbf{Attenuation} & Narrowing authority when passing it on (fewer fields, shorter life, one purpose). \\
    \textbf{Revocation} & Withdrawing authority, including anything derived from it. \\
    \textbf{Workload identity} & Proof of \emph{which software} is calling (SPIFFE), not for whom. \\
    \bottomrule
  \end{tabularx}
\end{frame}

\begin{frame}{The running scenario}
  \centering
  % The running scenario used in all three workshops (agentsim/world.py).
% Optional scale: \def\scenarioscale{0.8} before \input.
\providecommand{\scenarioscale}{1}
\begin{tikzpicture}[scale=\scenarioscale, transform shape, >=Stealth,
  ent/.style={draw=black!60, rounded corners=2pt, fill=white, minimum height=8mm, minimum width=19mm,
              align=center, font=\scriptsize, inner sep=2pt},
  org/.style={draw=#1, fill=#1!7, rounded corners=4pt, inner sep=5pt, thick},
  orglabel/.style={font=\tiny\bfseries, text=#1},
  hop/.style={->, thick, black!75},
  side/.style={->, thick, dashed, #1}]
  \node[ent] (alice) {Alice\\[-1pt]\tiny operator};
  \node[ent, right=5mm of alice] (orch) {orchestrator\\[-1pt]\tiny agent};
  \node[ent, right=12mm of orch] (xrf) {XRF analyzer\\[-1pt]\tiny agent};
  \node[ent, right=12mm of xrf] (reg) {DPP registry\\[-1pt]\tiny MCP server};
  \node[ent, below=5mm of reg, minimum width=16mm] (p1) {battery-001\\[-1pt]\tiny owned here};
  \node[ent, below=2mm of p1, minimum width=16mm] (p2) {battery-002\\[-1pt]\tiny othermaker};
  \node[ent, below=11mm of xrf, draw=orgbroker] (harv) {harvester\\[-1pt]\tiny broker agent};
  \node[ent, below=11mm of orch] (inc) {incident agent\\[-1pt]\tiny regulator};
  \begin{scope}[on background layer]
    \node[org=orgrecycler, fit=(alice)(orch)] (recbox) {};
    \node[org=orglab, fit=(xrf)] (labbox) {};
    \node[org=orgmaker, fit=(reg)(p1)(p2)] (mkbox) {};
    \node[org=orgbroker, fit=(harv), dashed] (brbox) {};
    \node[org=orgregulator, fit=(inc)] (rgbox) {};
  \end{scope}
  \node[orglabel=orgrecycler, above=0pt of recbox] {recycler};
  \node[orglabel=orglab, above=0pt of labbox] {lab};
  \node[orglabel=orgmaker, above=0pt of mkbox] {manufacturer};
  \node[orglabel=orgbroker, below=0pt of brbox] {broker (adversary)};
  \node[orglabel=orgregulator, below=0pt of rgbox] {regulator};
  \draw[hop] (alice) -- node[above, font=\tiny]{1} (orch);
  \draw[hop] (orch) -- node[above, font=\tiny]{2} (xrf);
  \draw[hop] (xrf) -- node[above, font=\tiny]{3} (reg);
  \draw[black!50] (reg) -- (p1);
  \draw[side=orgbroker] (harv.east) -- (harv.east -| mkbox.west);
\end{tikzpicture}

  \vspace{2mm}
  \begin{itemize}\small
    \item Alice asks the recycler's orchestrator to analyse battery-001; it hands the job to a third-party lab's XRF agent, which reads the passport from the manufacturer's registry.
    \item Two hops, three organizations. The broker's agent and the regulator's incident agent appear in Workshops 2 and 3.
    \item Everything is simulated in \texttt{agentsim}: mock authorization servers, MCP server, SPIFFE attestor, A2A cards. No LLM.
  \end{itemize}
\end{frame}

\begin{frame}{When is it really multi-agent?}
  \small
  \begin{tabularx}{\textwidth}{@{}l l X@{}}
    \toprule
    \textbf{Hop} & \textbf{Boundary} & \textbf{What crosses it} \\
    \midrule
    1 Alice $\to$ orchestrator & same organization & user authority into a workload \\
    2 orchestrator $\to$ XRF agent & recycler $\to$ lab & delegated authority to another owner's agent \\
    3 XRF agent $\to$ registry & lab $\to$ manufacturer & a request against data a third party owns \\
    \bottomrule
  \end{tabularx}
  \vspace{3mm}
  \begin{itemize}\small
    \item A single agent with tools would do if every tool were run by the same organization.
    \item Hop 2 is the real reason for multi-agent here: the lab owns the XRF capability and will not hand it over.
  \end{itemize}
  \prep{Add the classification of the DPP, reverse-logistics, and public-safety scenarios from my notes; cite R1.1 (Agent Index) only for claims it actually makes.}
\end{frame}

\section{What the protocols say}

\begin{frame}{MCP authorization in one slide}
  \small
  \begin{itemize}
    \item The MCP server is an OAuth resource server \ev{B}{mcp_authorization_2025}.
    \item It advertises its authorization server through protected-resource metadata \ev{B}{rfc_resourcemetadata_2025}; clients discover the server's endpoints \ev{B}{rfc_asmetadata_2018} and may register dynamically \ev{B}{rfc_dynamicregistration_2015}.
    \item Clients use PKCE \ev{B}{rfc_pkce_2015} and name the target with a resource indicator, so tokens carry an audience \ev{B}{rfc_resourceindicators_2020}.
    \item The server must check the audience and must not pass a received token through to upstream APIs.
  \end{itemize}
  \prep{Quote the exact MUST/SHOULD clauses with section numbers from notes R1.2; re-read the current spec revision before the workshop.}
\end{frame}

\begin{frame}{Demo E1.1: the flow, step by step}
  \centering\result{mcp_flow}
  \note{Run live: python workshops/w1-where-authority-leaks/experiments/e1_mcp_walkthrough.py. Pause on step 10: there is no standard step for a user from another organization. Step 11 succeeds and nobody can say for whom.}
\end{frame}

\begin{frame}{What the standards fix, and what they leave open}
  \centering\result{mcp_open}\par
  \vspace{2mm}
  {\small Governance-level critiques of the protocols: \ev{A}{governance_gaps_2026}\,\pre.}
\end{frame}

\begin{frame}{Identity is not authority}
  \begin{columns}[T]
    \column{0.55\textwidth}
    \result{cards}\par
    \vspace{3mm}
    \begin{itemize}\scriptsize
      \item A2A Agent Cards can be signed by the publishing organization \ev{B}{a2a_protocol_2026}.
      \item SPIFFE SVIDs prove which workload is calling \ev{B}{spiffe_standard_2026}. In the demo, both the lab's and the broker's agents authenticate.
    \end{itemize}
    \column{0.41\textwidth}
    \begin{exampleblock}{A verified card proves}\scriptsize
      the organization vouches for the description, and it was not altered.
    \end{exampleblock}
    \begin{alertblock}{It does not prove}\scriptsize
      which user the agent acts for, what it may access, or that its advertised skills are honest.
    \end{alertblock}
  \end{columns}
\end{frame}

\begin{frame}{Demo E1.2: passing tokens along}
  \centering\result{passthrough}\par
  \vspace{3mm}
  \begin{itemize}\small
    \item Audience checks stop a token from being used at the wrong server.
    \item They do not stop a thief who holds a token for the right server: a bearer token works for whoever presents it.
    \item Binding the token to the caller (here: SVID must match the token's actor) closes that gap. Standard options include DPoP \ev{B}{rfc_dpop_2023} and certificate-bound tokens \ev{B}{rfc_mtls_2020}.
  \end{itemize}
\end{frame}

\section{Delegation across organizations}

\begin{frame}{RFC 8693 token exchange}
  \begin{columns}[T]
    \column{0.48\textwidth}
    \textbf{Impersonation}
    \begin{itemize}\small
      \item The new token looks as if it had been issued to the subject.
      \item Downstream services cannot tell an agent was involved.
    \end{itemize}
    \column{0.48\textwidth}
    \textbf{Delegation}
    \begin{itemize}\small
      \item An \texttt{act} claim names the current actor; earlier actors nest inside it.
      \item For access decisions, only the top-level claims and the \emph{current} actor count. Prior actors are informational \ev{B}{rfc_tokenexchange_2020}.
    \end{itemize}
  \end{columns}
  \vspace{3mm}
  {\small Whether the new token's scope must be narrower is left to the authorization server's policy.}
  \prep{Check the section number for the "informational only" rule in notes R1.7.}
\end{frame}

\begin{frame}{Demo E1.3: the chain, delegation mode}
  \centering\result{chain_hops}\par
  \vspace{2mm}
  {\small Each hop's token names the next recipient as audience and the acting agent in \texttt{act}.}
\end{frame}

\begin{frame}{Four ways to issue the same chain}
  \centering\result{chain_modes}\par
  \vspace{2mm}
  \begin{itemize}\small
    \item \textbf{Impersonation:} the registry believes Alice is calling. The agents vanish from the record.
    \item \textbf{Scope copied:} the lab agent can read \emph{health}, which the task never needed.
    \item \textbf{No narrowing check:} the orchestrator asks for \emph{suppliers}, which Alice never had, and gets it.
    \item Narrowing is a property of one authorization server's configuration, not of the protocol.
  \end{itemize}
\end{frame}

\begin{frame}{Workload identity or delegation provenance?}
  \begin{columns}[T]
    \column{0.48\textwidth}
    \begin{exampleblock}{Workload identity is enough when}\small
      \begin{itemize}
        \item the agent acts for its own organization, on its own authority;
        \item no human principal needs to appear in the decision or the audit trail.
      \end{itemize}
    \end{exampleblock}
    \column{0.48\textwidth}
    \begin{alertblock}{Provenance is needed when}\small
      \begin{itemize}
        \item the agent acts for a person or another organization;
        \item authority must shrink, expire, or be revoked along the chain;
        \item someone must later reconstruct who authorized what.
      \end{itemize}
    \end{alertblock}
  \end{columns}
  \vspace{2mm}
  {\small Proposals that bind the two: \ev{A}{aip_delegation_2026}\,\pre, \ev{A}{hdp_provenance_2026}\,\pre; identity landscape: \ev{A}{aiidentity_standards_2026}\,\pre.}
  \prep{Summarize in one line each what AIP and HDP actually propose, from notes R1.10 and R1.11.}
\end{frame}

\section{Your turn}

\begin{frame}[fragile]{Hands-on: build the delegation chain (30 min)}
  \small
  \begin{enumerate}
    \item In pairs, fill in the activity sheet: for each hop, the issuer, subject, actor, audience, and scope.
    \item Copy your answers into \texttt{experiments/chain.json}.
    \item Run the structural verifier:
\begin{verbatim}
python e3_delegation_chain.py --check chain.json
\end{verbatim}
    \item Compare with the reference chain from Demo E1.3.
  \end{enumerate}
  \begin{block}{Questions to answer}\small
    Where does Alice's authority stop being visible? Which hop could amplify scope? Who could revoke hop 2?
  \end{block}
\end{frame}

\discussion{1}{Where should cross-organizational governance live: in the protocol, in a policy overlay, or in a registry?}
{MCP authorization \ev{B}{mcp_authorization_2025}; governance gaps in agent interoperability protocols \ev{A}{governance_gaps_2026}\,\pre. \prep{Prepare one argument for and one against each option.}}

\discussion{2}{When is workload identity enough, and when do we need delegation provenance?}
{RFC 8693 \ev{B}{rfc_tokenexchange_2020}; AIP \ev{A}{aip_delegation_2026}\,\pre; HDP \ev{A}{hdp_provenance_2026}\,\pre. \prep{Prepare the counter-case where provenance is overkill.}}

\begin{frame}{What carries into Workshop 2}
  \begin{enumerate}
    \item Authority leaks where the standards are silent: cross-organization users, scope per hop, revocation, audit.
    \item Signed identity (cards, SVIDs) says who is calling, not what they may do or for whom.
    \item RFC 8693 can carry the chain, but decisions may use only the current actor.
  \end{enumerate}
  \vspace{2mm}
  \textbf{Next:} can a policy model express ``the lab agent may read materials for Alice, for this task, until she revokes it''?
\end{frame}

\begin{frame}{How this workshop was prepared}
  \centering\result{prep_summary}
\end{frame}

\begin{frame}[allowframebreaks]{References}
  \printbibliography[heading=none]
\end{frame}

\end{document}
